\subsection*{一段安静时间没有带来新的证据}

设 Alice 在 09:00 用 \code{S1=(1,purchase.create.v1,E1,K1)} 购买 \code{paper-001}。本节初始余额为 10000 分，Q7 的成交金额为 3000 分，尚无购买。E1 的固定截止是 10:00。服务完成 P1/O1/D1，余额成为 7000，原结果中的 \code{balance\_after\_minor} 也为 7000。响应没有到达 Alice。她的设备随后失去网络，09:20 恢复，按已保存的 E1、K1 和原输入重试。

假设清理程序规定：完成记录保留十分钟，超过时长就删除 M 中的条目，入口看见 M 为空便按首次购买处理。09:10 的清理已经删除 M(S1)，09:20 的重试便生成 P2/O2/D2，余额变成 4000。客户端的等待没有改变原购买已经完成的事实，清理程序却删掉了入口用来识别这项事实的信息。

TTL 是按经过时长判定对象是否保留的规则。这个缩写在这里指记录的存活时间。它可以帮助选择何时运行清理，安全性则取决于清理以后入口怎样解释同一身份。RIFL 在讨论垃圾回收时也给出这种失败：客户端可能被分区或暂停得比记录的过期时间更久，恢复后再次发送已经完成的操作~\cite{rifl}。

\CTwelveFigure{ch12-original01}{逐笔 TTL 删除完成映射以后，有效 E1 下的旧 K1 被误当作首次购买}{fig:ch12-ttl}

把“客户端通常只重试几分钟”加入说明，尚未改变这个历史。网络可以保存一份迟到请求，进程可以在退出后重新启动，用户也可以从持久任务恢复原输入。09:20 的入口若继续允许 E1/K1 取得首次受理资格，就要找到一份足以判断 S1 已经受理的依据。完成记录只保留到 09:10，却让同一身份在 10:00 以前随时重新参加首次仲裁，两项时间规则便留下了五十分钟的缺口。

可以沿两条路径补上这段缺口。第一条路径保留已见 S1 的依据，先从响应副本中节省空间。09:10 可以删除热缓存，让持久档案继续重建完整原结果；Alice 尚未确认取得结果，原来的回放约定便保留到 10:00。只有客户端已经持久确认，或事先公布的回放期限已经结束，才删除最后一份可重建载荷，之后按已见标记拒绝重复首次执行。第二条路径把这一整段身份的首次受理资格也在删除前关闭，让入口以后根据关闭决定拒绝 E1/K1。前者保存逐笔依据，后者保存整段资格边界。清理程序先核对回放义务和后续判断的接替依据，再删除相应记录。

\CTwelveFigure{ch12-retention-boundary}{缓存副本、可重建结果与逐笔身份依据，分别按各自条件保留}{fig:ch12-retention-boundary}

\subsection*{把无法区分的两条历史摆在一起}

再构造两条独立历史，它们共同从余额 10000、空购买库和有效 E1 开始。两条历史的商品、报价、数量和金额完全相同。历史 A 中，09:00 到达的是 K1，购买完成，响应丢失。历史 B 中，09:00 到达的是另一项合法购买 K2，K1 的首次请求一直留在网络里。两条历史的第一笔购买都取得实际编号 P1/O1/D1，余额都成为 7000。

09:10，错误清理程序删除完成映射。09:20，入口收到同一份 K1 请求，使用同一观察过程：验证 E1 有效，按 S1 查询 M，得到空，再读取当前余额 7000。把这个观察写为 $\mathcal O$，它包含当前请求、会话资格、M(S1) 的空结果和余额。两条历史给出的 $\mathcal O$ 相同。

历史 A 的正确动作是恢复原购买，或在已经公布的回放结束规则下拒绝旧身份。再减一次 3000 会重复这项购买。历史 B 的 K1 尚未受理，而且还在 E1 的准入范围内，入口应允许它参加首次购买仲裁。此时可以再完成一项独立的 3000 分购买。相同 $\mathcal O$ 对应不同的正确动作，说明当前观察缺少“这份 S1 是否已经取得过绑定”的信息。

\CTwelveFigure{ch12-original02}{两条历史给出相同入口观察，历史 A 应保护旧效果，历史 B 可以首次受理}{fig:ch12-indistinguishable}

这个比较明确限定了入口读取什么。如果 P1 长期保存完整 S，历史 A 的 P1 属于 K1，历史 B 的 P1 属于 K2，业务数据能够区分两者。恢复入口实际查询这份身份索引，并在同一权威仲裁中使用它，就可以找回缺失的信息。仅仅说“订单还在”，而后续代码只查询已经清空的 M，无法利用这种差别。若身份关联也被一并删去，P/O/D 只剩相同的商品、金额与编号，两份业务状态就连这条恢复线索也失去了。

加大 TTL 可以让反例中的断网时间更长，改变不了信息缺口。另一种改进是让 09:20 的 K1 已经失去首次受理资格。此时两条历史都可返回过期，代价是历史 B 也无法再用这个旧身份首次购买。协议通过收紧准入范围，换取删除逐笔历史的空间。第三章登记的固定期限 E1，正是这样一条准入边界。

\subsection*{先说明要保留哪一种能力}

同一份记录可能承担三项工作。它让服务识别 S 已经受理，让客户端取回当时的完整结果，也让恢复者找到未完成支付。删除记录以前，分别判断这三项工作已经由什么信息接替。

完整原结果的要求很强。P1 成功时余额为 7000，后续 P2 成功后当前余额为 4000。订单查询即使能找到 O1，也未必包含 P1 当时的 7000、原拒绝原因或原响应的全部字段。只要接口继续承诺原结果回放，这些历史字段就有保存或可精确重建的位置。将一份结果从热存储移入可定位的持久档案，可以降低常用存储的占用，原结果的逻辑保留责任继续由档案承担。

防止重复首次受理所需的信息可以更少。服务只要可靠地认出 S 已退休，就可以拒绝再次执行。这里把退休用于描述回放能力已经按协议结束、旧身份也不会重新变成新操作的状态。它没有表达购买成功或失败。原结论即使是余额不足，也已经形成这项身份的终态，后来充值以后，原操作继续按这份拒绝结束。

未完成工作的责任另有期限。A 已提交 P1/C1/T1，平台结果未知时，删除响应缓存不会结束 T1。E1 到期以后，客户端购买入口可以返回过期，服务端继续沿 C1 查询原付款，保存证据并完成 R。把这项任务与 E1 的截止一起清空，会丢掉已经接受的工作。

\CTwelveFigure{ch12-original03}{原结果回放、旧身份拒绝与未完成恢复，依赖不同的保留依据}{fig:ch12-three-duties}

因此有限重试约定至少要说清两个时点：旧身份何时停止取得新的受理机会，原结果何时允许停止回放。第三章公布的 \code{retry\_until} 给出了固定期限，普通客户端在此之前可按原规则恢复。接下来增加的客户端确认，是调用方主动释放某些已取得结果的回放权利。没有发送这种确认的客户端，按原期限保留结果。
\FloatBarrier
